\section{Modelo mecánico: robot, carga y contacto}
\label{sec:modelo-fisico}

Las ecuaciones de \cref{sec:dinamica} bastan para derivar la cota de esfuerzo,
pero no exhiben de dónde salen las restricciones ni qué distingue realmente los
dos modos de transporte. Esta sección desarrolla el modelo completo.

\subsection{El AMR: formulación lagrangiana con restricción no holónoma}

Sea $q_i=(x_i,y_i,\theta_i)$ el vector de coordenadas generalizadas. La rueda no
desliza lateralmente, lo que impone la restricción de Pfaff
\begin{equation}
  \dot x_i\sin\theta_i-\dot y_i\cos\theta_i=0
  \qquad\Longleftrightarrow\qquad
  A(q_i)\,\dot q_i=0,
  \quad A(q_i)=\bigl(\sin\theta_i,\ -\cos\theta_i,\ 0\bigr).
  \label{eq:pfaff}
\end{equation}
Esta restricción no es integrable: no proviene de ninguna relación entre las
coordenadas, solo entre sus velocidades. De ahí el término \emph{no holónoma} y
de ahí que el AMR pueda alcanzar cualquier pose y no pueda hacerlo en cualquier
dirección instantánea.

Aplicando Euler--Lagrange con multiplicador $\lambda_i$ asociado a
\eqref{eq:pfaff},
\begin{equation}
  M(q_i)\ddot q_i+C(q_i,\dot q_i)\dot q_i+F(\dot q_i)
  \;=\;B(q_i)\,\tau_i+A(q_i)^{\!\top}\lambda_i,
  \label{eq:lagrange}
\end{equation}
con $M$ simétrica definida positiva, $C$ los términos centrípetos y de Coriolis,
$F$ la fricción, $B$ la matriz de entrada y $\tau_i=(\tau_{iL},\tau_{iR})$ los
pares de rueda.

\paragraph{Reducción y eliminación del multiplicador.}
Sea $S(q_i)$ una base del núcleo de $A$, es decir $A(q_i)S(q_i)=0$:
\begin{equation}
  S(q_i)=
  \begin{pmatrix}\cos\theta_i&0\\ \sin\theta_i&0\\ 0&1\end{pmatrix},
  \qquad
  \dot q_i=S(q_i)\,\eta_i,\quad \eta_i=(v_i,\omega_i).
  \label{eq:reduccion}
\end{equation}
Premultiplicando \eqref{eq:lagrange} por $S^{\!\top}$ desaparece el término
$S^{\!\top}A^{\!\top}\lambda_i=0$ y queda el modelo reducido
\begin{equation}
  M_1\dot\eta_i+C_1(q_i,\dot q_i)\,\eta_i+F_1=B_1\tau_i,
  \qquad
  M_1=S^{\!\top}MS,\ \ B_1=S^{\!\top}B,\ \
  C_1=S^{\!\top}M\dot S+S^{\!\top}CS,
  \label{eq:reducido}
\end{equation}
que es la forma usada en \eqref{eq:ruedas}. El número de columnas de $S$ es
exactamente el grado de movilidad $\delta_m$ de \cref{def:movilidad}: la
reducción lagrangiana y la clasificación cinemática son el mismo objeto visto
desde dos lados.

\begin{observacion}[El multiplicador es la fuerza que las ruedas deben dar]
\label{obs:multiplicador-friccion}
Eliminar $\lambda_i$ del modelo no lo elimina de la física. $A^{\!\top}\lambda_i$
es la fuerza lateral de restricción, y quien la ejerce es el rozamiento
neumático--suelo. La restricción \eqref{eq:pfaff} solo se cumple mientras
\[
  |\lambda_i|\ \le\ \mu_{\mathrm{sd}}\,m_ig ,
\]
y cuando se excede, la rueda desliza y el modelo reducido
\eqref{eq:reducido} deja de describir el sistema. La no holonomía no es una
propiedad geométrica incondicional: es una restricción que se sostiene mientras
haya adherencia. Es el mismo límite que aparece en \cref{pr:cota}, ahora con su
origen identificado.
\end{observacion}

\subsection{Un solo disco de fricción por rueda}

\Cref{obs:multiplicador-friccion} acota el multiplicador lateral por
$\mu_{\mathrm{sd}}m_ig$ como si la sujeción lateral dispusiera de un presupuesto
propio. No lo tiene: la fuerza de contacto de cada rueda vive en \emph{un} disco
de Coulomb, y la tracción longitudinal y la reacción lateral lo comparten
\cite{mayorga2026m1}.

\begin{proposicion}[Paso de fricción de máxima disipación]
\label{pr:disco-friccion}
Con geometría y normales congeladas, sean $v^{*}$ las velocidades tras las fuerzas
no incluidas en el subpaso, $M\succ0$ la inercia y $A$ el mapa de deslizamientos
de ambas ruedas. El impulso de fricción es
\[
  j^\star\in\arg\min_{j\in\Dcal}\ \tfrac12\,j^{\!\top}AM^{-1}A^{\!\top}j+(Av^{*})^{\!\top}j,
  \qquad
  \Dcal=\bigl\{j:\ \lVert j_s\rVert_2\le\mu_sN_s\Delta t,\ s\in\{L,R\}\bigr\},
\]
y $v^{+}=v^{*}+M^{-1}A^{\!\top}j^\star$. Las dos componentes de cada rueda comparten
el disco: no hay presupuestos longitudinal y lateral independientes.
\end{proposicion}

\begin{observacion}[Lo que cuesta empujar fuerte: sujeción lateral]
\label{obs:disco-compartido}
Sobre el subproblema de una rueda con $R=\mu N\Delta t=0{,}180$~N\,s y una demanda
lateral fija, la sujeción lateral disponible cae de $0{,}180$ a $0{,}176$,
$0{,}164$ y $0{,}145$ al pedir tracción longitudinal creciente, con el impulso
siempre en la frontera del disco. Es decir: \emph{cuanto más empuja un AMR, menos
capacidad le queda para no derrapar}. Esto refina la lectura de
\cref{obs:multiplicador-friccion}: la cota lateral no es una constante sino lo
que sobra del disco tras la tracción, y la parálisis por adherencia de
\cref{th:degradacion} llega antes de lo que sugiere el presupuesto lateral aislado.

Hay además una razón de modelado para este núcleo. La regularización viscosa
$F_y=-C_y v_\perp$ es disipativa pero da $F_y=0$ cuando $v_\perp=0$: no puede
representar una reacción lateral \emph{estática}, que es exactamente lo que
sostiene a una coalición acoplada en reposo. El programa anterior sí la produce
—impulso lateral $0{,}180$ con velocidad transversal nula— y por eso es el que
se ejecuta en la planta de \cref{sec:ensayos}. No se afirma que sea una
implementación completa de los métodos impulso–momento en que se inspira.
\end{observacion}

\figDisco

\subsection{Tres propiedades que hacen tratable el modelo}

Escribir \eqref{eq:lagrange} no basta: lo que permite controlarlo son tres
propiedades estructurales de todo sistema de Euler--Lagrange
\cite{cai2022cooperative}, que conviene enunciar porque cada una sostiene una
pieza distinta de este trabajo.

\begin{proposicion}[Propiedades estructurales]
\label{pr:propiedades-el}
Para el modelo \eqref{eq:lagrange} con $C$ construida con los símbolos de
Christoffel:
\begin{enumerate}[label=(\roman*),leftmargin=2.2em,topsep=0.2em,itemsep=0.15em]
  \item $M(q)\succeq k_m I$ con $k_m>0$, luego $M(q)^{-1}$ está acotada
  uniformemente;
  \item $\dot M(q)-2C(q,\dot q)$ es antisimétrica, y en consecuencia
  $z^{\!\top}\bigl(\dot M-2C\bigr)z=0$ para todo $z$;
  \item la dinámica es lineal en los parámetros: existe una matriz de regresión
  \emph{conocida} $Y$ tal que
  \[
    M(q)x+C(q,\dot q)y+G(q)\;=\;Y(q,\dot q,x,y)\,\Theta
  \]
  para todo $x,y$, con $\Theta$ el vector constante de parámetros inciertos.
\end{enumerate}
\end{proposicion}

\begin{observacion}[Qué sostiene cada una]
\label{obs:propiedades-uso}
La primera acota $\lVert M^{-1}\rVert$ y es lo que permite despejar
aceleraciones sin que las cotas de \cref{pr:cota} degeneren.

La segunda es la propiedad de pasividad, y es la razón de que la energía sirva
como función de Lyapunov: el término de Coriolis no realiza trabajo neto, de modo
que al derivar $\tfrac12\dot e^{\!\top}M\dot e$ ese término se cancela
exactamente en lugar de tener que acotarse. Toda la construcción disipativa de
\cref{th:smallgain} descansa sobre esta cancelación en la capa mecánica.
Comprobación: sobre un sistema de dos grados de libertad con configuraciones
aleatorias, $\lVert(\dot M-2C)+(\dot M-2C)^{\!\top}\rVert_\infty$ resulta del
orden de $10^{-7}$, que es el ruido de la diferenciación numérica, y
$\dot q^{\!\top}(\dot M-2C)\dot q\approx4\times10^{-9}$.

La tercera es la más consecuente para \cref{obs:adaptacion}. Allí se declaró que
la incertidumbre paramétrica queda abierta; ahora puede decirse algo más
preciso: esa incertidumbre entra \emph{linealmente}, con regresor conocido, y esa
linealidad es exactamente lo que hace posible estimarla en línea. La
comprobación numérica da igualdad a precisión de máquina. Lo que separa
\cref{obs:adaptacion} de un estimador operativo no es, por tanto, una dificultad
estructural, sino la excitación persistente.
\end{observacion}

\subsection{Lo que la adaptación cierra, y lo que deja igual}

\Cref{obs:propiedades-uso} deja la incertidumbre paramétrica a un estimador con
excitación persistente. La identidad que ese estimador necesita puede escribirse
con precisión, y con ella se ve qué cierra y qué no \cite{mayorga2026m1}.

\begin{proposicion}[Balance adaptativo condicionado]
\label{pr:balance-adaptativo}
Dentro de una rama regular con entradas realizables, sea
$M(q)\dot v+C(q,\dot q)v+g(q)=u+d$, $v_r$ una referencia diferenciable y
$s=v-v_r$. Supóngase la parametrización
$M\dot v_r+Cv_r+g=Y(q,\dot q,v_r,\dot v_r)\,\theta$ de
\cref{pr:propiedades-el}(iii) y $\dot M-2C$ antisimétrica. Con el mando y la
adaptación
\[
  u=Y\hat\theta-Ks,\qquad \dot{\hat\theta}=-\Gamma Y^{\!\top}s,\qquad K\succ0,\ \Gamma\succ0,
\]
y un residual total $r$ que recoge perturbación y diferencia entre entrada
solicitada y realizada, el almacenamiento
$V=\tfrac12 s^{\!\top}Ms+\tfrac12\tilde\theta^{\!\top}\Gamma^{-1}\tilde\theta$
satisface $\dot V=-s^{\!\top}Ks+s^{\!\top}r$; y si $k=\lambda_{\min}(K)$,
$\dot V\le-\tfrac{k}{2}\lVert s\rVert^{2}+\lVert r\rVert^{2}/(2k)$.
\end{proposicion}

\begin{proof}
Al derivar, $s^{\!\top}(\dot M/2-C)s=0$ por antisimetría; los términos
$s^{\!\top}Y\tilde\theta$ y $\tilde\theta^{\!\top}\Gamma^{-1}\dot{\tilde\theta}$ se
cancelan por la ley de adaptación. Quedan el amortiguamiento y la potencia del
residual.
\end{proof}

\begin{observacion}[Comprobado, con sus tres límites]
\label{obs:adaptacion-num}
Sobre el sistema de dos grados de libertad de \cref{obs:propiedades-uso} con
parámetros iniciales errados hasta un $50\%$: con $r=0$ y sin excitación, $s$
cae a $2{,}7\times10^{-3}$ mientras $\lVert\tilde\theta\rVert$ solo baja de
$1{,}42$ a $1{,}07$; con excitación, $s$ cae a $4{,}6\times10^{-3}$ y
$\lVert\tilde\theta\rVert$ de $0{,}59$ a $0{,}055$; con $r\neq0$, $s$ queda
acotado en $3{,}4\times10^{-2}$ sin anularse. La identidad $\dot V=-s^{\!\top}Ks$
se verifica por diferencias con error medio $3\times10^{-5}$ en el caso sin
excitación.

Tres límites que la derivación declara y que se mantienen aquí. Primero, la
desigualdad no es una prueba de estabilidad entrada--estado del vector completo
$(s,\tilde\theta)$: el término negativo no domina el error paramétrico, y
$\tilde\theta\to0$ no se deduce sin excitación —el ensayo lo muestra—. Proyección
de estimaciones, fugas o hipótesis de excitación exigen un análisis adicional y no
se añaden sin modificar la prueba. Segundo, la adaptación reduce el error de
realización de una continuación, pero el juego y el filtro de barrera deben
seguir contabilizando $r$: una saturación de par o una pérdida de contacto no
desaparecen por actualizar $\hat\theta$. Tercero, si la planta está subactuada
$u$ no puede tratarse como fuerza arbitraria, que es exactamente la situación de
\cref{pr:fuera-eje} cuando el punto controlado está sobre el eje.
\end{observacion}

\begin{observacion}[Parámetros propios y parámetros compartidos]
\label{obs:parametros-compartidos}
Hay además una separación que la arquitectura debe respetar. Las masas e inercias
de los AMR son \emph{propias}: no debe imponerse consenso entre las $m_i$ de AMR
distintos, y cada uno estima lo suyo. La masa y la inercia de la carga son
\emph{compartidas}: su estimación cooperativa debe distinguir qué medidas son
independientes y cuáles son copias de una misma observación, que es el mismo
problema de correlación que \cref{pr:interseccion} resuelve de forma conservadora
para las poses. Estimar la carga promediando ingenuamente las estimaciones de
sus portadores vuelve a contar la misma evidencia, y por \cref{th:asimetria} no
hay un consenso exacto disponible para ese objeto.
\end{observacion}

\subsection{El punto controlado y por qué conviene sacarlo del eje}

Un AMR industrial no acepta pares: acepta \emph{referencias de velocidad}
$(u^{r},\omega^{r})$, que su electrónica interna sigue. El modelo utilizable en
control es entonces el de \cite{gorrostieta2019applications}, donde
\eqref{eq:reducido} se completa con la dinámica de seguimiento
\begin{equation}
  \begin{pmatrix}\dot u\\ \dot\omega\end{pmatrix}
  =\begin{pmatrix}
     \tfrac{\theta_3}{\theta_1}\omega^{2}-\tfrac{\theta_4}{\theta_1}u\\[2pt]
    -\tfrac{\theta_5}{\theta_2}u\,\omega-\tfrac{\theta_6}{\theta_2}\omega
   \end{pmatrix}
  +\begin{pmatrix}\tfrac{1}{\theta_1}&0\\ 0&\tfrac{1}{\theta_2}\end{pmatrix}
   \begin{pmatrix}u^{r}\\ \omega^{r}\end{pmatrix}
  +\begin{pmatrix}\delta_u\\ \delta_\omega\end{pmatrix},
  \label{eq:vel-dinamica}
\end{equation}
con $\theta=(\theta_1,\dots,\theta_6)$ parámetros identificables
experimentalmente y $\delta$ la incertidumbre paramétrica. Las cotas de
\cref{pr:cota} se trasladan a $(u^{r},\omega^{r})$ por esta relación.

La observación decisiva es sobre \emph{qué punto} se controla. Sea $h$ el punto
situado a distancia $a$ del centro del eje a lo largo del rumbo —el pusher, la
pinza, el punto de contacto—. Su velocidad es
\begin{equation}
  \dot h=M(\psi,a)\begin{pmatrix}u\\ \omega\end{pmatrix},
  \qquad
  M(\psi,a)=\begin{pmatrix}\cos\psi&-a\sin\psi\\ \sin\psi&a\cos\psi\end{pmatrix},
  \qquad \det M=a .
  \label{eq:fuera-eje}
\end{equation}

\begin{proposicion}[La no holonomía no restringe al punto de contacto]
\label{pr:fuera-eje}
Para $a\neq0$, $M(\psi,a)$ es invertible con valores singulares $1$ y $|a|$,
luego $\sigma_{\min}(M)=\min(1,|a|)$. En consecuencia el punto $h$ puede recibir
\emph{cualquier} velocidad planar, y la velocidad angular requerida para una
velocidad lateral unitaria crece como $1/a$.
\end{proposicion}

\begin{proof}
$M^{\!\top}M=\operatorname{diag}(1,a^{2})$, de donde los valores singulares. Para
$a\neq0$ el sistema $\dot h=Mv$ tiene solución única
$v=M^{-1}\dot h$, y la componente angular de $M^{-1}$ escala como $1/a$.
\end{proof}

\begin{observacion}[La ley de degradación en la geometría del chasis]
\label{obs:fuera-eje-degradacion}
La comprobación numérica da $\sigma_{\min}=\min(1,a)$ exacto, y para $1$~m/s
lateral se requieren $15{,}3$~rad/s con $a=0{,}05$~m, $3{,}82$ con $a=0{,}2$ y
$0{,}76$ con $a=1$: el producto $|\omega|\,a$ es constante. El offset $a$ es,
por tanto, un parámetro de \emph{autoridad} con la misma ley $1/a$ de
\cref{th:degradacion}, ahora no en la geometría del contacto sino en la del
chasis. Un pusher pegado al eje es cinemáticamente casi singular; alejarlo
compra maniobrabilidad de forma exactamente cuantificada. Es una consecuencia de
diseño mecánico que sale del mismo teorema.
\end{observacion}

\subsection{Funciones transversales: lo que cuesta sacar el punto del eje}

\Cref{pr:fuera-eje} muestra que un punto adelantado del eje admite cualquier
velocidad planar. La construcción general de ese hecho son las \emph{funciones
transversales}: coordenadas auxiliares invertibles sobre las que se diseña el
controlador, con un error geométrico ajustable respecto de la pose física. Lo que
esa construcción añade a este trabajo es una cota explícita de ese error, que es
lo que las barreras necesitan \cite{mayorga2026m1}.

\begin{proposicion}[Del error transversal a una huella segura]
\label{pr:huella-transversal}
Sean $(\bar p,\bar\phi)$ las coordenadas físicas y $(p,\phi)$ las auxiliares,
$p=\bar p+R(\phi)(f_1,f_2)^{\!\top}$, $\phi=\bar\phi-f_3$, con
$f_3=\epsilon_2\cos\xi$, $0<\epsilon_2<\pi/2$,
$f_1=\epsilon_1\sin\xi\,\sin f_3/f_3$ y $f_2=\epsilon_1\sin\xi\,(1-\cos f_3)/f_3$.
Entonces $\lVert p-\bar p\rVert\le\epsilon_1$ y $|\bar\phi-\phi|\le\epsilon_2$, y
para un cuerpo compacto contenido en un disco de radio $R_b$ la distancia de
Hausdorff entre sus huellas física y auxiliar está acotada por
\[
  \delta_b=\epsilon_1+2R_b\sin(\epsilon_2/2),
\]
y si la pose auxiliar tiene además error de seguimiento $(\varepsilon_p,\varepsilon_\phi)$
respecto de una pose nominal,
$\delta_b\le\epsilon_1+\varepsilon_p+2R_b\sin\bigl(\min\{\pi,\epsilon_2+\varepsilon_\phi\}/2\bigr)$.
\end{proposicion}

\begin{observacion}[Comprobación y uso]
\label{obs:huella-num}
Sobre trescientas poses aleatorias por caso, el desplazamiento máximo de todo
punto material del cuerpo entre las dos huellas queda en el $77$–$79\%$ de
$\delta_b$ para $(\epsilon_1,\epsilon_2)$ de $(0{,}05;\,0{,}2)$ a
$(0{,}2;\,1{,}0)$ con $R_b=0{,}6$: la cota se respeta con margen. El uso es
directo: cuando el controlador actúa sobre la pose auxiliar, las barreras de
\cref{pr:hocbf} deben tensarse en $\delta_b$ además del tensado por incertidumbre
de \cref{co:tensado}. Son dos inflaciones distintas de la huella —una por
estimación, otra por diseño del controlador— y se suman.
\end{observacion}

\begin{observacion}[El precio en mando, con número]
\label{obs:mando-grande}
Reducir el error geométrico no tiene coste de actuación nulo, y el caso más
simple lo fija. Para el punto adelantado $z=p+\ell\,e(\theta)$ se tiene
$\omega=(Se)^{\!\top}\dot z/\ell$: con una velocidad lateral auxiliar de
$0{,}2$~m/s y $|\omega|\le1{,}5$~rad/s hace falta $\ell\ge0{,}133$~m, y usar
$\ell=0{,}05$~m exigiría $4$~rad/s. Es la ley $1/a$ de
\cref{obs:fuera-eje-degradacion} convertida en regla de dimensionamiento del
pusher: $\ell_{\min}=v_{\mathrm{lat}}/\omega_{\max}$. Un \emph{offset} pequeño
compra precisión geométrica y la paga en demanda angular, y el control virtual
adicional no es un motor nuevo: la no holonomía no desaparece, se desplaza a la
saturación de $\omega$.
\end{observacion}

\figAMR

\subsection{La carga y cómo la afectan los contactos}

La carga es un cuerpo rígido planar de pose $\xi_\ell\in SE(2)$ y \emph{twist}
$\zeta_\ell=(V_x,V_y,\Omega)$. Su dinámica en el marco del cuerpo es
\begin{equation}
  M_\ell\dot\zeta_\ell+C_\ell(\zeta_\ell)\zeta_\ell+D_\ell\zeta_\ell
  \;=\;\sum_{i\in\Ccal_\ell}G_{i}\,f_i\;+\;w^{\mathrm{ext}}_\ell,
  \qquad
  M_\ell=\operatorname{diag}(m_\ell,m_\ell,I_\ell),
  \label{eq:carga-dinamica}
\end{equation}
donde $D_\ell$ recoge la resistencia al avance —fricción de arrastre contra el
suelo, que para una carga apoyada es $\mu_s m_\ell g$ y no un término viscoso—,
$w^{\mathrm{ext}}_\ell$ las perturbaciones, y $G_i$ el mapa que lleva la fuerza
de contacto del AMR $i$ al \emph{wrench} resultante:
\begin{equation}
  G_i=\begin{pmatrix} R(\psi_\ell)^{\!\top} & 0\\[2pt] (r_i)^{\perp\top} & 1\end{pmatrix},
  \qquad
  G_i f_i=\begin{pmatrix} f_i^{\mathrm{lin}}\\ r_i^{\perp}\!\cdot f_i^{\mathrm{lin}}+m_i^{c}\end{pmatrix},
  \label{eq:mapa-contacto}
\end{equation}
con $r_i$ el vector del centro de masa al punto de contacto y $m^c_i$ el par de
contacto, nulo si el contacto es puntual sin fricción de giro.

\begin{observacion}[Por qué el brazo importa tanto como la fuerza]
\label{obs:brazo}
\Cref{eq:mapa-contacto} muestra que la contribución al par es
$r_i^{\perp}\!\cdot f_i^{\mathrm{lin}}$: solo cuenta la componente de la fuerza
perpendicular al brazo. Un AMR que empuja radialmente hacia el centro de masa
aporta fuerza y \emph{cero} par, por grande que sea su empuje. Esa es la razón
geométrica de que la cardinalidad no baste (\cref{th:noequiv}) y de que el rango
tampoco (\cref{th:span}): la capacidad de una coalición depende de dónde toca,
no solo de cuánto empuja.
\end{observacion}

\subsection{Modo \textsf{cargo}: acoplamiento bilateral}

En \textsf{cargo} el AMR se acopla a un puesto mediante un elemento de contacto
capaz de transmitir tracción y compresión —una pinza, un enganche o un apoyo con
retención—. El contacto es \emph{bilateral}: la fuerza puede tener cualquier
signo en la dirección normal.

Sea $\rho_i$ el anclaje material en el marco de la carga y
$\epsilon_i=p_i-\xi_\ell\oplus\rho_i$ el error de acoplamiento. Un modelo de
acoplamiento elástico con amortiguamiento da
\begin{equation}
  f_i=K_{\mathrm{pad}}\,\epsilon_i-D_{\mathrm{pad}}\,u_i,
  \qquad
  u_i=H(r_i)\zeta_\ell-\bigl(v_ie(\theta_i)+\omega_iJR(\theta_i)a_i\bigr),
  \label{eq:pad}
\end{equation}
con $u_i$ la velocidad relativa en el contacto. La condición de adherencia
$\|f_i\|\le\mu_{c,i}N_i$ delimita el régimen en que \eqref{eq:pad} es válido.

Las propiedades del modo se siguen de la bilateralidad:
\begin{enumerate}[leftmargin=1.9em,topsep=0.2em,itemsep=0.2em]
  \item el conjunto realizable es un \emph{subespacio} local, no un cono, luego
  la caracterización de \cref{th:span} se relaja: basta rango completo si las
  cotas no saturan;
  \item aparecen fuerzas internas $\bm\lambda^{0}\in\ker G_\Ccal$ que aprietan
  sin mover (\cref{th:rigidizacion}), con el compromiso de
  \cref{co:agarre};
  \item la formación queda sujeta a la restricción de rigidez
  $A_{\mathrm{kin}}\zeta_\ell=0$ de \cref{pr:rigidez}, que puede bloquear
  maniobras que cada AMR podría ejecutar por separado.
\end{enumerate}

\begin{proposicion}[Identidad energética del \emph{pad} isotrópico]
\label{pr:pad-energia}
Con el acoplamiento \eqref{eq:pad} escrito como $f_i=k_i\epsilon_i-d_iu_i$ y
$\dot\epsilon_i=-u_i+\Omega S\epsilon_i$, con fuerzas opuestas aplicadas en el
mismo punto y $k_i,d_i>0$,
\[
  f_i^{\!\top}u_i+\frac{d}{dt}\Bigl(\tfrac12 k_i\lVert\epsilon_i\rVert^{2}\Bigr)
  \;=\;-d_i\lVert u_i\rVert^{2}.
\]
\end{proposicion}

\begin{proof}
Sustituir ambas ecuaciones cancela $k_i\epsilon_i^{\!\top}u_i$; el término de
rotación es nulo porque $\epsilon_i^{\!\top}S\epsilon_i=0$.
\end{proof}

\begin{observacion}[El acoplamiento es pasivo, y eso es lo que se compone]
\label{obs:pad-pasivo}
La identidad se verifica a $4{,}5\times10^{-17}$ sobre trayectorias aleatorias, y
dice que el \emph{pad} es un elemento pasivo con almacenamiento
$\tfrac12k_i\lVert\epsilon_i\rVert^{2}$ y disipación $d_i\lVert u_i\rVert^{2}$: es
exactamente la clase de bloque que \cref{th:smallgain} compone, y por eso el
acoplamiento elástico entra en la red disipativa sin hipótesis adicional. Su
límite lo declara la planta: la ley no puede conservar su anclaje como si
siguiera adherida cuando sale del cono de fricción, y la política posterior a una
pérdida general del apoyo requiere un modelo de deslizamiento y readquisición que
no está aquí \cite{mayorga2026m1}.
\end{observacion}

\begin{proposicion}[Levantar una ruta de carga hasta cada rueda]
\label{pr:ruta-ruedas}
Para el apoyo $i$ en $r_i=R(\psi_\ell)\rho_i$, la velocidad del punto material es
$u_i=V_\ell+\Omega S r_i$ y su aceleración
$\dot u_i=\dot V_\ell+\dot\Omega Sr_i-\Omega^{2}r_i$. Para seguir la curva sin
deslizamiento lateral, el rumbo del AMR debe girar a
\[
  \omega_i=\frac{(Su_i)^{\!\top}\dot u_i}{\lVert u_i\rVert^{2}},\qquad u_i\neq0,
\]
y las velocidades de rueda son $\omega_{w,iL/R}=(v_i\mp b_i\omega_i)/r_{w,i}$ con
$v_i=e_i^{\!\top}u_i$. El espacio de \emph{twists} de carga compatibles con la
rodadura, antes de aplicar cotas de velocidad, tiene dimensión
$3-\operatorname{rank}A_{\mathrm{kin}}$.
\end{proposicion}

\begin{observacion}[Lo que esto añade a la rigidez, y su singularidad]
\label{obs:ruta-ruedas}
La expresión de $\omega_i$ reproduce la derivada numérica del ángulo del vector
velocidad con error $1{,}6\times10^{-5}$ sobre una trayectoria curva de la
carga. Es la conexión que faltaba entre \cref{pr:rigidez} y \cref{pr:cota}: una
ruta de la carga no impone solo compatibilidad de \emph{twist}, impone una
velocidad angular y unas velocidades de rueda concretas a cada AMR, y son estas
las que deben caber en sus cotas. Dos precisiones. La dimensión
$3-\operatorname{rank}A_{\mathrm{kin}}$ vale \emph{antes} de cotas y con el
contacto sobre el eje, que es exactamente el régimen que
\cref{obs:rigidez-alcance} acota; y un \emph{swivel} no hace omnidireccional al
AMR: permite reorientar el chasis en reposo, no elimina la restricción lateral
instantánea. La fórmula de $\omega_i$ es singular en $u_i=0$ —apoyo momentáneamente
parado—, y ahí no se usa: se planifica la continuación de orientación, que es una
de las guardas de \cref{def:automata}.
\end{observacion}

\subsection{Modo \textsf{caging}: contacto unilateral y confinamiento}

En \textsf{caging} no hay agarre. Los AMR rodean la carga y la empujan; el
contacto es \emph{unilateral} y obedece la complementariedad
\eqref{eq:complementariedad} con la ley de Coulomb \eqref{eq:coulomb}. Por
\cref{pr:notiron} ningún contacto tira: solo comprime.

La diferencia esencial con \textsf{cargo} no es la ley de contacto sino
\emph{qué se garantiza}. En \textsf{cargo} se controla la pose de la carga; en
\textsf{caging} se garantiza que la carga no escapa, y su pose dentro de la
jaula queda parcialmente libre.

\begin{definicion}[Jaula]
\label{def:jaula}
La coalición $\Ccal$ en configuración $q$ \emph{enjaula} a la carga si la
componente conexa del espacio libre de configuración de la carga que la contiene
es acotada.
\end{definicion}

\Cref{def:jaula} es topológica y no dice cómo comprobarla. El siguiente
certificado es geométrico, conservador y calculable \cite{mayorga2026m1}.

\begin{teorema}[Cerco por polígono de centros]
\label{th:cerco}
Supóngase que la carga contiene, para toda orientación, un disco de radio $r_L$
centrado en su referencia $P_L$, y que cada AMR contiene un disco rígido ocupado
de radio $r_i$ centrado en $p_i$. Sean $p_1,\dots,p_m$ los vértices de un
polígono simple cerrado, de área no nula, que contiene estrictamente a $P_L$. Si
para cada índice cíclico
\[
  \lVert p_{i+1}-p_i\rVert\ \le\ (r_i+r_L)+(r_{i+1}+r_L),
  \qquad
  \lVert P_L-p_i\rVert\ >\ r_i+r_L ,
\]
entonces la referencia de la carga no puede salir del interior sin colisión entre
discos ocupados.
\end{teorema}

\begin{proof}
Cada disco de radio $r_i+r_L$ en torno a $p_i$ es región prohibida para $P_L$.
Sobre el segmento $[p_i,p_{i+1}]$ los intervalos cubiertos desde sus extremos se
solapan por la primera desigualdad, luego todo el segmento pertenece a la unión
prohibida, y con él la frontera poligonal completa. Por el teorema de la curva de
Jordan, una curva continua que abandone el interior cruza esa frontera y por
tanto una región prohibida. La segunda desigualdad garantiza que el estado
inicial es libre.
\end{proof}

\begin{corolario}[Cerco móvil condicionado]
\label{co:cerco-movil}
La pertenencia al interior se conserva para un cerco en movimiento si centros y
carga evolucionan de forma continua, el polígono permanece simple y no
degenerado, y sus segmentos siguen cubiertos en todo instante.
\end{corolario}

\begin{observacion}[Suficiente, no necesario, y comprobado]
\label{obs:cerco-num}
Con seis AMR de radio ocupado $0{,}30$~m y carga con disco inscrito $0{,}35$~m,
un cerco de radio $1{,}10$~m cumple las desigualdades y un muestreo de cuatro
mil trayectorias rectas de escape no encuentra ninguna libre; con radio
$1{,}60$~m las desigualdades fallan y el escape se encuentra. Tres salvedades
que el propio desarrollo declara y que se mantienen aquí: un círculo exterior
que solo encierre las ruedas no sirve, porque el disco ocupado debe ser rígido;
el criterio es suficiente y no necesario —una carga rectangular puede quedar
confinada aunque su disco inscrito no lo reconozca—; y bajo incertidumbre,
apretar las distancias entre vecinos no certifica simplicidad robusta del
polígono, de modo que la rutina rechaza la afirmación robusta cuando no dispone
de esa prueba. Se incorpora precisamente para no llamar \emph{caging} a
cualquier empuje.
\end{observacion}

\begin{observacion}[Tres condiciones, tres exigencias]
\label{obs:jerarquia-contacto}
Conviene no confundir tres nociones de la teoría del agarre:
\emph{cierre de forma}, que inmoviliza la carga por geometría sin fricción;
\emph{cierre de fuerza}, que permite generar un \emph{wrench} arbitrario con
fricción; y \emph{enjaulado}, que solo impide el escape. La primera es la más
exigente y la última la más débil, y esa debilidad es su ventaja: enjaular exige
menos AMR y menos precisión de contacto. \Cref{pr:caging} da una condición
suficiente por curva cerrada, y es una condición topológica —de componente
conexa— no de fuerzas. Nada de la biblioteca de robótica móvil consultada
formula esta jerarquía; procede de la literatura de manipulación
\cite{rimon1996caging,bicchi1995closure}.
\end{observacion}

\begin{observacion}[Las dos resiliencias]
\label{obs:dos-resiliencias}
De \cref{obs:jerarquia-contacto} se sigue una diferencia operativa importante
ante el fallo de un miembro. En \textsf{cargo}, perder un AMR reduce el conjunto
realizable y puede invalidar el certificado: la resiliencia es una propiedad
\emph{dual}, verificable con el margen de \cref{th:robusto} sobre la coalición
reducida. En \textsf{caging}, perder un AMR solo importa si abre la jaula: la
resiliencia es una propiedad \emph{topológica}, verificable sobre componentes
conexas. Son certificados de naturaleza distinta y ninguno implica el otro; un
modo híbrido —jaula más un mínimo de puestos de \emph{wrench}— exige comprobar
los dos.
\end{observacion}

\figCarga

\subsection{Los límites de movimiento y cómo se componen}

Cinco restricciones acotan lo que la coalición puede hacer, y actúan a la vez:

\begin{table}[!ht]
\centering\small
\caption{Los límites del movimiento y su origen.}
\label{tab:limites}
\begin{tabularx}{\textwidth}{@{}p{3.5cm}p{4.6cm}X@{}}
\toprule
Límite & Expresión & Origen \\
\midrule
Par de actuador & $|\tau_{i\{L,R\}}|\le\bar\tau_i$ & motor \\
Adherencia de tracción & $|\lambda_i|\le\mu_{\mathrm{sd}}m_ig$
  & validez de \eqref{eq:pfaff} (\cref{obs:multiplicador-friccion}) \\
Esfuerzo de contacto & $0\le\lambda\le\bar\lambda_i$ con
  $\bar\lambda_i$ de \cref{pr:cota} & mínimo de par y adherencia \\
Curvatura y velocidad & $v_i\le v^{\max}/(1+|\kappa_i|b_i/2)$
  & $\delta_s=0$ (\cref{le:curvatura}) \\
Frenado & $v_i\le\sqrt{2a^{\max}_i(h_{\min}-d_{\min})}$
  & factibilidad del QP (\cref{le:astop}) \\
\bottomrule
\end{tabularx}
\end{table}

\begin{observacion}[El límite activo cambia con la situación]
\label{obs:limite-activo}
Los cinco no son redundantes y ninguno domina siempre. En un pasillo despejado
manda la curvatura; junto a un obstáculo, el frenado; con carga pesada, la
adherencia; en maniobra de giro con contacto rígido, la compatibilidad
$A_{\mathrm{kin}}$. Reportar cuál está activo en cada instante es diagnóstico
barato y explica detenciones que de otro modo parecen arbitrarias: cada una es
un límite distinto saturando.
\end{observacion}

\begin{teorema}[El grado relativo incluye las ruedas]
\label{th:grado-relativo-rueda}
En la reducción longitudinal con rueda dinámica y contacto flexible,
\[
  F_g=c_g(r_w\omega-v),\quad f=k_t(x-P)+d_t(v-V),\quad
  M\dot V=f-b_LV,\quad m\dot v=F_g-f-b_av,\quad J_w\dot\omega=\tau-r_wF_g-b_w\omega ,
\]
con $c_g,r_w,m,M,J_w,d_t>0$, la posición $P$ de la carga tiene grado relativo
cuatro respecto del par $\tau$, con coeficiente $d_tc_gr_w/(MmJ_w)$ en
$P^{(4)}$; la posición del AMR tiene grado relativo tres, con coeficiente
$c_gr_w/(mJ_w)$; y si $d_t=0$ y $k_t>0$, el de la carga sube a cinco con
coeficiente $k_tc_gr_w/(MmJ_w)$.
\end{teorema}

\begin{proof}
$\dot P=V$ y $\ddot P=(f-b_LV)/M$ dependen del estado; $\dddot P$ contiene
$\dot f$, que contiene $\dot v$ y con él $F_g$, que depende de $\omega$ pero no
del par instantáneo; derivar una vez más introduce $\dot\omega$ y con él
$c_gr_w\tau/J_w$, multiplicado por $d_t/(Mm)$. Si $d_t=0$ esa vía desaparece
y hace falta otra derivada, a través de $k_t\dot v$. Las tres cadenas se
comprobaron con derivación simbólica sobre el modelo declarado.
\end{proof}

\begin{observacion}[Barrera de cuarto orden y sensibilidad $T^4$ del precio]
\label{obs:grado-cuatro}
Dos consecuencias. La barrera de \cref{obs:rigidez-alcance} sobre la
posición de la carga debe derivarse hasta cuarto orden —con ganancias iguales,
$h^{(4)}+4\alpha h^{(3)}+6\alpha^2\ddot h+4\alpha^3\dot h+\alpha^4h\ge0$, con
las cuatro condiciones auxiliares iniciales no negativas— y no reutilizar una
barrera de aceleración como si el par entrase en primer orden. Y un precio de
ocupación espacial que actúe sobre la posición predicha tiene sensibilidad de
orden $T^4$ al par en esta reducción, $\tfrac{c\,T^4}{24}B_P\,\tau$, de modo
que la constante de la ventana de \cref{th:N-ventana} depende de la física
actuada y no puede heredarse de una planta sin rueda. Saturación, cambio de
contacto y anulación del coeficiente direccional degradan la autoridad; el
teorema es del modelo ideal.
\end{observacion}
